We need to consider the asymptotic behavior of $\expect{\post}{\phi(\theta, \xvec_s)}$,
which can be written as the ratio of the integrals
%
\begin{align*}
    %
\expect{\post}{\phi(\theta, \xvec_s)} ={}&
\frac{\int \phi(\theta, \xvec_s)
        \exp\left(N \likhat(\theta)\right)
        d \theta}
  {\int \exp\left(N\likhat(\theta)\right) d \theta}.
 %
\end{align*}
%
Recall that $\thetahat := \mathrm{argmax}_\theta \likhat(\theta)$ and
define $\thetatil(t) := t (\theta - \thetahat) + \thetahat$.  We will
employ the following Taylor series expansion:
\begin{align}
\likhat(\theta) ={}&
    \likhat(\thetahat) +
    \likhatk{1}(\thetahat)(\theta - \thetahat) +
    \frac{1}{2}\likhatk{2}(\thetahat)(\theta - \thetahat)^2 +
\nonumber
\\&
    \frac{1}{6}\likhatk{3}(\thetahat)(\theta - \thetahat)^3 +
    \tresid{4}{\likhat, \thetahat, \theta} (\theta - \thetahat)^4,
    \eqlabel{taylor_expansion}
\end{align}
%
recalling the definition of the integral form of the Taylor series residual
$\tresid{4}{\likhat, \thetahat, \theta}$ given in \defref{taylor_residual}.

Note that $\likhat(\thetahat)$ is constant in $\theta$, and that
$\likhatk{1}(\thetahat) = 0$.  Plugging in and changing the domain of
integration to $\tau$, we see
%
\begin{align*}
\expect{\post}{\phi(\theta, \xvec_s)}
={}&
\frac{\int \phi(\theta, \xvec_s)
       \exp\left(
            N \likhat(\theta) - N \likhat(\thetahat)
        \right)
       d \theta}
     {\int
     \exp\left(
        N \likhat(\theta) - N \likhat(\thetahat)
     \right) d \theta}
\\={}&
\frac{\int \phi(\theta, \xvec_s)
       \exp\left(
            N \likhat(\theta) - N \likhat(\thetahat)
        \right)
       \sqrt{N} d \tau}
     {\int
     \exp\left(
        N \likhat(\theta) - N \likhat(\thetahat)
     \right) \sqrt{N} d \tau}.
% =\\{}&
% \frac{\int \phi(\theta, \xvec_s)
%        \exp\left(
%             \frac{1}{2} \likhatk{2}(\thetahat)\tau^2 +
%             N^{-1/2} \frac{1}{6}\likhatk{3}(\thetahat) \tau^3 +
%             N^{-1} \tresid{4}{\likhat, \thetahat, \theta} \tau^4
%         \right)
%        d \tau}
%      {\int
%      \exp\left(
%          \frac{1}{2} \likhatk{2}(\thetahat)\tau^2 +
%          N^{-1/2} \frac{1}{6}\likhatk{3}(\thetahat) \tau^3 +
%          N^{-1} \tresid{4}{\likhat, \thetahat, \theta} \tau^4
%      \right) d \tau}.
\end{align*}
%
Since the factor of $\sqrt{N}$ cancels in the preceding display, it will suffice
to consider the following generic integral.
%
\begin{defn}\deflabel{bclt_integral}
%
For a target function $\phi$, define the integral functional
%
\begin{align}
I(\phi) :={}&
\int \phi(\theta, \xvec_s)
   \exp\left(
       N \likhat(\theta) - N \likhat(\thetahat)
   \right)
   d \tau \\
={}&
\int \phi(\theta, \xvec_s) \times
\\&\quad
   \exp\left(
       \frac{1}{2} \likhatk{2}(\thetahat)\tau^2 +
       N^{-1/2} \frac{1}{6}\likhatk{3}(\thetahat) \tau^3 +
       N^{-1} \tresid{4}{\likhat, \thetahat, \theta} \tau^4
   \right)
   d \tau.
   \eqlabel{integral_funcational}
\end{align}
%
\end{defn}
%
Under \defref{bclt_integral}, $\expect{\post}{\phi(\theta, \xvec_s)} = I(\phi) /
I(1)$.  We will now consider the behavior of $I(\phi)$ for a
generic third-order BCLT-okay $\phi$,
returning to consider the behavior of the ratio in \secref{bclt_region_comb}.
